\chapter*{INTRODUCTION GÉNÉRALE}
\addcontentsline{toc}{chapter}{INTRODUCTION GÉNÉRALE}

L'imagerie médicale a révolutionné la médecine moderne en offrant la capacité de visualiser l'intérieur du corps humain de manière non invasive, permettant ainsi un diagnostic précoce, un suivi thérapeutique précis et une planification chirurgicale optimale. Parmi les différentes modalités d'imagerie, la \textbf{tomodensitométrie (TDM)}, communément appelée scanner, occupe une place prépondérante grâce à sa capacité à fournir des images anatomiques détaillées avec une excellente résolution spatiale. Cependant, cette technologie repose sur l'utilisation de rayons X, dont l'exposition répétée soulève des préoccupations croissantes en termes de risques radio-induits pour les patients, en particulier dans les populations sensibles comme les enfants ou les patients nécessitant des examens de suivi fréquents. Cette problématique de réduction de dose d'irradiation tout en maintenant une qualité d'image diagnostique constitue l'un des défis majeurs de la recherche contemporaine en imagerie médicale.\vspace{7pt}\\
La reconstruction d'image tomographique consiste à estimer une carte des coefficients d'atténuation des tissus à partir des projections acquises par le scanner. Historiquement, la méthode de référence est la \textbf{rétroprojection filtrée (FBP)}, une technique analytique rapide et robuste qui repose sur le théorème de la coupe centrale et la transformée de Fourier. La \textbf{FBP} a démontré son efficacité pendant des décennies et reste largement utilisée en routine clinique en raison de sa simplicité et de sa rapidité d'exécution. Néanmoins, cette méthode présente des limitations fondamentales qui deviennent particulièrement critiques dans un contexte de réduction de dose. En effet, lorsque le nombre de projections diminue ou que le bruit augmente, la \textbf{FBP} produit des images dégradées par des artefacts de stries prononcés et une amplification significative du bruit, rendant l'interprétation clinique difficile, voire impossible. Ces limitations sont inhérentes à sa nature analytique qui ne permet pas d'intégrer efficacement des connaissances \textit{a priori} sur la structure des images médicales.\vspace{7pt}\\
Parallèlement, la théorie du \textbf{Compressed Sensing (CS)}, introduite au début des années 2000 par les travaux fondateurs de Candès, Romberg, Tao et Donoho, a ouvert une voie nouvelle pour la reconstruction d'images à partir de données sous-échantillonnées. Le principe fondamental du \textbf{CS} repose sur deux piliers : la parcimonie, qui stipule qu'un signal peut être représenté de manière concise dans un domaine de transformation approprié, et l'incohérence, qui garantit que l'acquisition sous-échantillonnée préserve l'information essentielle. Appliqué à la \textbf{TDM}, le \textbf{CS} offre la perspective séduisante de reconstruire des images de haute qualité à partir d'un nombre réduit de projections, permettant ainsi une réduction significative de la dose d'irradiation. Cette approche a suscité un intérêt considérable dans la communauté scientifique et a conduit au développement de nombreuses méthodes de reconstruction itératives intégrant une régularisation parcimonieuse.\vspace{7pt}\\
La formulation mathématique des problèmes de reconstruction par \textbf{CS} conduit typiquement à un problème d'optimisation composite combinant un terme d'attache aux données, qui mesure la fidélité de la reconstruction par rapport aux mesures acquises, et un terme de régularisation parcimonieuse, qui encourage une représentation simple de l'image dans un domaine transformé. Ce type de problème, bien que conceptuellement élégant, pose des défis numériques importants en raison de sa nature non lisse et de la grande dimensionnalité des données. La résolution efficace de ces problèmes d'optimisation nécessite des algorithmes sophistiqués capables de gérer la structure composite tout en garantissant une convergence rapide et stable. Parmi ces algorithmes, la méthode des multiplicateurs à direction alternée (\textbf{ADMM}) s'est imposée comme un outil de choix, grâce à sa capacité à décomposer le problème original en sous-problèmes plus simples traités alternativement, permettant ainsi de séparer la gestion de l'attache aux données de celle de la régularisation parcimonieuse.\vspace{7pt}\\
Dans ce contexte, le choix de la représentation parcimonieuse est tout aussi crucial que le choix de l'algorithme d'optimisation. Les transformées en ondelettes occupent une place privilégiée pour la représentation des images naturelles et médicales, car elles offrent une analyse multi-résolution qui capture efficacement à la fois les structures lisses et les discontinuités que sont les contours. Parmi la vaste famille des ondelettes, les ondelettes biorthogonales présentent des propriétés particulièrement attractives pour la reconstruction d'images. Contrairement aux ondelettes orthogonales, les ondelettes biorthogonales peuvent être symétriques, une propriété essentielle pour éviter les distorsions de phase qui décalent ou déforment les contours lors de la reconstruction. La famille \textbf{Bior4.4}, en particulier, avec ses quatre moments de vanishing pour l'analyse et la synthèse, offre un compromis optimal entre la représentation des zones homogènes et la préservation des détails anatomiques fins.\vspace{7pt}\\
Ce travail de thèse s'inscrit dans cette problématique de réduction de dose en \textbf{TDM} par des méthodes de reconstruction avancées. L'objectif principal est de développer et valider une approche de reconstruction tomographique combinant l'algorithme \textbf{ADMM} avec la transformée en ondelettes biorthogonales \textbf{Bior4.4}, et de démontrer sa supériorité par rapport aux méthodes conventionnelles dans un contexte de sous-échantillonnage angulaire. L'innovation de cette recherche ne réside pas dans l'invention d'outils mathématiques à partir de rien, mais dans la démonstration rigoureuse que le couplage spécifique entre l'\textbf{ADMM} et la \textbf{Bior4.4} constitue une solution particulièrement performante et robuste, comme en attestent les résultats quantitatifs obtenus avec un gain de plus de \SI{12}{dB} en \textbf{PSNR} par rapport à la \textbf{FBP} et un \textbf{SSIM} de 0.963 témoignant d'une préservation structurelle quasi-parfaite.\vspace{7pt}\\
La structure de ce manuscrit reflète la progression logique de notre démarche scientifique. Le premier chapitre (\ref{chapitre1}) pose les fondements de la reconstruction d'image en présentant les principales modalités d'imagerie médicale, avec un accent particulier sur la \textbf{TDM}, la transformation de Radon qui la sous-tend, et les limites des méthodes analytiques conventionnelles comme la \textbf{FBP}. Ce chapitre introduit également les concepts fondamentaux du \textbf{CS} et leur application à la reconstruction tomographique. Le deuxième chapitre (\ref{chapitre2}) est consacré aux outils mathématiques utilisés dans ce travail, avec une présentation détaillée de la transformée de Fourier, de la convolution, de la transformée de Radon et de sa discrétisation, aboutissant à la formulation linéaire du problème de reconstruction. Le troisième chapitre (\ref{chapitre3}) constitue le cœur théorique de notre contribution, avec une présentation approfondie du problème inverse en tomodensitométrie, de sa régularisation parcimonieuse, de l'algorithme \textbf{ADMM} et de ses fondements, ainsi que des transformées en ondelettes biorthogonales et de leurs propriétés. Enfin, le quatrième chapitre (\ref{chapitre4}) présente les aspects pratiques de notre travail, incluant la simulation des données, l'implémentation de l'algorithme, l'analyse détaillée des résultats quantitatifs et de la convergence, et une discussion approfondie sur l'innovation de notre approche et ses perspectives pour la réduction de dose en imagerie médicale.